\section{Introduction}\label{sec:intro}

Whether a minimum wage helps the workers it targets depends on how many of them it actually reaches. In developing countries a large share of employment lies outside the reach of labor regulation, so a statutory wage floor can raise pay for covered workers, displace some of them into the uncovered informal sector, or act only as a reference point for wages that formal firms would have paid in any case. Which of these forces dominates is an empirical question, and its answer turns on how binding the floor is, how well it is enforced, and how freely workers move between the formal and informal margins. I study these questions for Peru, a country that pairs a minimum wage high relative to median pay with one of the highest rates of labor informality in Latin America.

I exploit the increase in Peru's national minimum wage (the \emph{Remuneración Mínima Vital}, or RMV) from 930 to 1,025 soles, enacted by Supreme Decree 003-2022-TR and effective on 1 May 2022. This 10.2 percent nominal increase was the first change to the RMV in four years, and it occurred as the labor market was emerging from the pandemic. A central institutional fact shapes the analysis: Peru sets a single national minimum wage, with no regional or sectoral variation in the general regime. There is therefore no cross-sectional variation in the level of the policy and no staggered adoption to exploit. As \citet{jimenez_2023} notes for the 2016 reform, credible identification must instead come from differential \emph{exposure} to a policy that changes for everyone at the same time.

My source of differential exposure is education. In the pre-reform quarters the median monthly-equivalent wage of a low-education private wage earner was 950 soles, almost exactly the old minimum of 930, and 57 percent of low-education wage earners were paid below the \emph{new} floor of 1,025 soles on the eve of the reform; a ten percent increase in that floor is therefore highly binding for this group. High-education workers, by contrast, have a median wage well above the minimum and are less likely to be bound by it, although a substantial minority remains exposed through informal or own-account work. This gap in exposure, rather than a claim of zero exposure for the comparison group, motivates a difference-in-differences design that contrasts low-education workers (those with at most complete secondary schooling, about two thirds of the working-age population) with high-education workers (those with any tertiary schooling). The design applies a long tradition of measuring minimum wage effects from variation in the ``bite'' of the policy across groups or regions \citep{card_1992, lee_1999, machin_manning_rahman_2003} to a setting with a single national reform and rich, quarterly household data.

I use the public-use quarterly employment modules of the Encuesta Nacional de Hogares (ENAHO) from 2021Q1 through 2024Q4, which bracket the reform with five pre-reform and ten clean post-reform quarters; the transition quarter 2022Q2, in which the reform took effect mid-quarter, is excluded from the baseline and restored in a robustness check. The data are repeated cross-sections, and their quarterly frequency lets me estimate a dynamic event study, test for differential pre-trends by education, and trace the effect for more than two years after the reform. I examine a broad set of outcomes: real hourly wages and monthly earnings, the probability of employment and of labor force participation, the incidence of formal versus informal employment, hours of work, and the split of employment between wage work and self-employment. Throughout I weight by the survey weights and cluster inference by department, and because treatment is assigned at a coarse level I complement conventional cluster-robust standard errors with the small-sample correction of \citet{pustejovsky_tipton_2018} and with randomization inference.

Three findings emerge. First, and most robustly, the reform shifted the \emph{composition} of low-education employment away from formal wage work. The probability that a low-education worker holds formal employment falls by 4.6 percentage points, about a fifth of the group's 23 percent baseline formality rate, and self-employment rises by 3.6 points. These estimates have pre-reform formality trends that pass their joint test, with a small residual drift that I quantify and adjust for; they appear immediately after the reform and persist, are larger in the departments where the new floor cuts deeper into the wage distribution, and are confirmed by an independent continuous-exposure design that relates outcomes to the pre-reform regional Kaitz index and supports design-based randomization inference. Linked persons in the ENAHO Panel 2020--2024 reveal the flow behind the stock: incumbent formal workers were retained, while the rate at which informal low-educated workers entered formal jobs fell by nearly half, so the reallocation operated through the creation of formal matches rather than their destruction. Second, there is no evidence that the reform lifted the group's pay, but not because the floor is a dead letter. Within the formal sector, where enforcement operates, the earnings spike at the old minimum migrates almost completely to the new one, and the share of formal wage earners paid below the new floor falls; among informal wage earners that share \emph{rises}, and because only 13 percent of employed low-education workers are formal, the two responses cancel: the aggregate below-floor share does not move, no decile of the wage distribution shows a gain, and the estimated effect on the average log real hourly wage is negative and insignificant. The no-lift conclusion rests on the compliance and distributional evidence, whose diagnostics are clean; the average wage regression itself carries a rejected pre-trend test, driven by a single early-2021 quarter, and I am explicit about what it can and cannot rule out. Third, the employment rate of low-education workers falls by 2.5 percentage points, but this margin is the least robust: it fails in-time placebo tests, has a clear differential pandemic-recovery pre-trend, and does not scale with regional exposure, so I do not give it a causal interpretation. Throughout I bound conclusions with the sensitivity analysis of \citet{rambachan_roth_2023} and a battery of specification and placebo checks. I then replicate the whole design on the January 2025 reform, which raised the RMV to 1,130 soles in a calmer, post-pandemic window: no wage response appears there either, and neither does the reallocation, which I interpret as evidence that the margin along which the minimum wage binds depends on the state of the labor market.

The contribution is a combination rather than a new tool, and I state it without overclaiming. To my knowledge this is the first evaluation of the 2022 Peruvian minimum wage reform and the first study of any Peruvian reform using the post-pandemic quarterly ENAHO; the existing micro evidence stops at the 2016 reform \citep{jimenez_2023} or at the reforms of the early 2000s \citep{jaramillo_lopez_2006, cespedes_2005}. It is also the first to make heterogeneity by education the central object of a worker-level event study of the Peruvian minimum wage, where prior work cuts heterogeneity by income range or age or constructs a department-level dose of exposure. It is, to my knowledge, also the first to use the linked ENAHO Panel to measure worker-level formality transitions around any Peruvian minimum wage reform, which lets me separate the retention of incumbents from the entry of new workers into formal jobs. And it places the informal margin at the center of the analysis, showing that adjustment runs through the composition of employment rather than through pay. I claim no novelty for the identification idea, which is the standard exposure or bite design, nor for the econometric toolkit; the value added is the pairing of a new reform and a policy-relevant post-pandemic period with an education-based worker-level design and a careful treatment of informality and of the threats to identification that such a design entails.

The rest of the paper proceeds as follows. Section~\ref{sec:institutional} describes the Peruvian minimum wage and the reform, and Section~\ref{sec:literature} places the paper in the literature. Section~\ref{sec:data} presents the data and the construction of the key variables, and Section~\ref{sec:strategy} sets out the empirical strategy and the identifying assumptions. Section~\ref{sec:results} reports the main results, Section~\ref{sec:robustness} the robustness and inference checks, and Section~\ref{sec:heterogeneity} the heterogeneity and distributional analysis. Section~\ref{sec:discussion} discusses mechanisms, external validity, and limitations, and Section~\ref{sec:conclusion} concludes.
